\documentclass[11pt,a4paper]{article}
\usepackage[T1]{fontenc}
\usepackage[utf8]{inputenc}
\usepackage[brazilian]{babel}
\usepackage[margin=1in]{geometry}
\usepackage{setspace}
\usepackage{enumitem}
\usepackage{hyperref}
\usepackage{xurl}
\usepackage{parskip}

\setstretch{1.15}
\setlength{\parindent}{0pt}
\emergencystretch=2em
\setlist[enumerate]{leftmargin=2em}
\hypersetup{colorlinks=true,linkcolor=black,urlcolor=blue,citecolor=black}

\begin{document}

\begin{center}
{\large Instituto Brasileiro de Ensino, Desenvolvimento e Pesquisa}\par
Mestrado em Administração Pública -- Ciência de Dados e Inteligência Artificial\par
Turma: 2005.2\par
Disciplina: Inteligência Artificial Generativa\par
Professor: Marcelo Pita\par
Alunos: Analécia Rorato; Fabricio Fernandes Santana; Paulo Ceser Siqueira Neves

\vspace{1.5em}
{\LARGE \textbf{Chatbot legislativo com RAG para respostas auditáveis sobre discursos da 56\textordfeminine{} Legislatura do Senado Federal}}\par

\vspace{0.75em}
Repositório: \url{https://github.com/fabriciosantana/mcdia/tree/main/05-iag/4-project}
\end{center}

\section*{Resumo}

Este trabalho investiga a aplicação da abordagem retrieval-augmented generation (RAG) no desenvolvimento de um chatbot legislativo voltado à consulta de discursos da 56\textordfeminine{} Legislatura do Senado Federal (2019--2023). Parte-se do problema da dificuldade de acesso e recuperação de informação sobre os discursos parlamentares por meio de consultas em linguagem natural, especialmente quando se exige precisão factual, contextualização e rastreabilidade das informações. A metodologia consiste na construção de um pipeline de processamento que inclui extração e preparação de dados, segmentação textual (chunking), geração de embeddings, indexação vetorial e implementação de fluxo de consulta baseado em LLM. A solução foi desenvolvida em ambiente reprodutível, utilizando Open WebUI, ChromaDB, Ollama, API da OpenAI e scripts próprios para ingestão e avaliação. Os resultados indicam que o sistema responde adequadamente a perguntas factuais e temáticas, especialmente quando bem delimitadas, com base em trechos documentais relevantes. Observou-se, contudo, que consultas amplas ou multifocais reduzem a precisão das respostas, evidenciando a influência do chunking, da formulação da consulta e do desenho do prompt. Conclui-se que a abordagem RAG é viável e adequada para recuperação de discursos, promovendo ganhos em transparência, rastreabilidade, auditabilidade e acesso à informação legislativa.

\textbf{Palavras-chave:} RAG; inteligência artificial; chatbot legislativo; recuperação de informação; Senado Federal; LLM.

\section{Introdução}

\subsection{Contextualização}

O avanço recente dos grandes modelos de linguagem (large language models -- LLMs) ampliou significativamente o interesse por soluções capazes de responder perguntas em linguagem natural (MINAEE et al., 2024). No setor público, esse potencial mostra-se particularmente relevante em contextos caracterizados por acervos extensos, dinâmicos e heterogêneos, tais como legislação, pareceres, notas técnicas, discursos parlamentares e documentos administrativos (INTER-PARLIAMENTARY UNION, 2024). Nesses ambientes, o desafio não se restringe ao armazenamento da informação, mas abrange, sobretudo, sua recuperação qualificada, contextualizada e verificável (WESTMINSTER FOUNDATION FOR DEMOCRACY, 2023).

Apesar de seu potencial, o uso direto de LLMs apresenta limitações amplamente discutidas na literatura, entre as quais se destacam a alucinação, a desatualização do conhecimento usado no treinamento e a baixa rastreabilidade das evidências utilizadas na produção das respostas (GAO et al., 2023). Tais limitações assumem maior gravidade em contextos institucionais nos quais se exige transparência, verificabilidade, controle das fontes e responsabilização sobre os conteúdos produzidos (WESTMINSTER FOUNDATION FOR DEMOCRACY, 2023).

Nesse cenário, a abordagem retrieval-augmented generation (RAG) apresenta-se como alternativa tecnicamente promissora, pois, em vez de depender exclusivamente do conhecimento internalizado no modelo, sistemas RAG recuperam trechos relevantes de uma base documental externa e os incorporam ao contexto da geração (GAO et al., 2023). Com isso, tende-se a elevar a precisão factual, a atualidade e a auditabilidade das respostas, ao mesmo tempo em que se favorece a explicitação das fontes que as sustentam (GAO et al., 2023).

No contexto do Parlamento brasileiro, a organização do conhecimento legislativo e a ampliação do acesso qualificado à informação constituem desafios relevantes de pesquisa aplicada, como observa Martellote Viola et al. (2025) ao defender que o emprego de inteligência artificial nesse domínio deve ser acompanhado de atenção à estrutura informacional, à mediação documental e à explicabilidade dos resultados.

É nesse horizonte que se insere o presente trabalho, ao propor o desenvolvimento de um chatbot legislativo baseado em RAG, voltado à consulta de discursos da 56\textordfeminine{} Legislatura do Senado Federal.

\subsection{Motivação}

A administração pública contemporânea opera sob crescente pressão por transparência, rastreabilidade, eficiência e capacidade de resposta (UNITED NATIONS, 2015). No ambiente parlamentar, a consulta a discursos e pronunciamentos mostra-se relevante para diversas finalidades, tais como análise temática, comparação de posicionamentos, apoio técnico a gabinetes, elaboração de relatórios, produção de subsídios e atendimento ao cidadão (SKUBIC et al., 2024). Entretanto, a simples disponibilização pública desses documentos não garante, por si só, acesso eficiente ao conteúdo neles contido (INTERNATIONAL IDEA, 2023).

Em grandes acervos discursivos, mecanismos tradicionais de busca, predominantemente baseados em correspondência lexical, tendem a ser insuficientes quando o usuário formula perguntas complexas, comparativas, interpretativas ou orientadas à síntese (GUPTA et al., 2024). Além disso, a recuperação da informação frequentemente depende de conhecimento prévio sobre vocabulário, datas, autores ou a estrutura da base documental. Como consequência, a informação permanece formalmente acessível, mas substantivamente difícil de explorar.

Paralelamente, soluções baseadas exclusivamente em LLMs generalistas podem produzir respostas plausíveis e linguisticamente bem estruturadas, porém desprovidas de fundamentação verificável (ANH-HOANG et al., 2025; OPENAI, 2025). Em domínios institucionais, tal comportamento é particularmente problemático, pois compromete a confiança, a governança da informação e a responsabilização sobre o uso da tecnologia (HIGHBERG, 2025). Deene (2025) argumenta que sistemas RAG podem mitigar parte desses riscos ao deslocar a base do conhecimento para documentos controlados, atualizáveis e auditáveis. Na mesma direção, Said (2025) enfatiza que, em cenários de produção, sistemas RAG requerem práticas adicionais de versionamento, observabilidade e avaliação contínua, de modo a assegurar qualidade e consistência ao longo do tempo.

Dessa forma, a motivação deste trabalho decorre da necessidade de investigar uma solução tecnicamente viável e institucionalmente adequada para ampliar o acesso qualificado a discursos parlamentares, combinando busca semântica, geração textual e mecanismos de auditabilidade.

\subsection{Problema de pesquisa}

Diante desse contexto, formula-se o seguinte problema de pesquisa: como estruturar uma solução baseada em RAG capaz de permitir a consulta, em linguagem natural, a discursos da 56\textordfeminine{} Legislatura do Senado Federal, produzindo respostas úteis, fundamentadas em evidências recuperadas e compatíveis com exigências de auditabilidade em contexto institucional?

A formulação desse problema decorre da percepção de que o acesso documental, embora formalmente assegurado, ainda enfrenta limitações relevantes quando se exige recuperação semântica, síntese orientada ao usuário e rastreabilidade das fontes utilizadas na resposta.

\subsection{Objetivos}

O objetivo geral deste trabalho consiste em desenvolver e demonstrar uma solução de chatbot legislativo baseada em RAG para consulta a discursos parlamentares da 56\textordfeminine{} Legislatura do Senado Federal.

Como objetivos específicos, pretende-se:
\begin{enumerate}
  \item estruturar uma base de conhecimento consultável a partir dos discursos parlamentares;
  \item possibilitar a formulação de perguntas em linguagem natural com respostas fundamentadas em trechos recuperados;
  \item avaliar empiricamente a qualidade inicial dos processos de recuperação e geração;
  \item documentar uma arquitetura reprodutível, modular e passível de evolução para cenários institucionais.
\end{enumerate}

\subsection{Justificativa}

A relevância deste estudo decorre de sua contribuição em dois planos complementares. No plano prático, a pesquisa investiga uma forma de ampliar o acesso qualificado a informações parlamentares públicas, contribuindo para transparência, memória institucional e apoio a atividades de gabinete, pesquisa e controle social. No plano teórico-aplicado, o trabalho examina a aplicação de técnicas contemporâneas de recuperação e geração em um domínio documental complexo, marcado por exigências elevadas de explicabilidade, governança e rastreabilidade.

Além disso, o estudo se justifica pela aderência do problema a desafios concretos da administração pública digital. Em vez de tratar a inteligência artificial como mecanismo autônomo de produção de respostas, propõe-se abordagem em que a geração textual é condicionada por fontes documentais públicas, controláveis e auditáveis. Essa orientação é coerente com demandas institucionais por uso responsável de IA em contextos sensíveis.

\section{Fundamentação teórica}

\subsection{Arquiteturas de LLM relacionadas ao problema}

As arquiteturas mais diretamente relacionadas a este trabalho são, de um lado, os LLMs generativos baseados em transformers (VASWANI et al., 2017) e, de outro, os modelos de embeddings utilizados para recuperação semântica (NIE et al., 2024). No primeiro caso, o LLM atua como componente de síntese textual, recebendo uma pergunta e um conjunto de trechos recuperados para compor uma resposta. No segundo, modelos de embedding convertem documentos e consultas em vetores de alta dimensionalidade, tornando possível o cálculo de similaridade semântica em bancos vetoriais (GAO et al., 2023).

Em sistemas RAG, tais componentes operam de maneira complementar. O modelo de embedding sustenta a etapa de recuperação, enquanto o modelo generativo realiza a etapa de composição da resposta a partir do contexto recuperado (GAO et al., 2023). Essa separação funcional é relevante porque favorece modularidade arquitetural e especialização de componentes, o que, por sua vez, amplia a capacidade de evolução incremental da solução (GAO et al., 2024).

Segundo Gao et al. (2023), sistemas RAG podem ser classificados em abordagens ingênuas, avançadas e modulares. Para o problema investigado neste trabalho, a arquitetura modular mostra-se especialmente adequada, pois permite substituir, de forma relativamente independente, extratores documentais, estratégias de segmentação, modelos de embedding, bancos vetoriais, reranqueadores e modelos gerativos. Em contextos institucionais, essa característica é particularmente desejável, uma vez que restrições de infraestrutura, custos operacionais, requisitos de segurança e políticas de governança podem demandar alterações graduais da pilha tecnológica.

\subsection{Técnicas relacionadas}

As principais técnicas adotadas neste projeto incluem extração documental, chunking, geração de embeddings semânticos, recuperação vetorial, reranqueamento, seleção contextual e prompting orientado a evidências (SAID, 2025; MAHARJAN et al., 2026). A extração documental corresponde à conversão e preparação do texto para ingestão na base de conhecimento. O chunking refere-se à segmentação dos documentos em unidades menores, adequadas à indexação vetorial e ao uso posterior como contexto de geração. Os embeddings semânticos representam vetorialmente documentos e consultas, enquanto a recuperação vetorial identifica os trechos semanticamente mais relevantes. O reranqueamento e a seleção contextual, por sua vez, delimitam o subconjunto de trechos efetivamente enviado ao modelo generativo. Por fim, o prompting orientado a evidências consiste na formulação de instruções explícitas para que o modelo privilegie o contexto recuperado, evite inferências especulativas e apresente referências de modo transparente.

Dentre essas técnicas, o chunking ocupa posição central. Singh (2025) argumenta que estratégias inadequadas de segmentação podem comprometer significativamente a recuperação semântica, seja por fragmentarem indevidamente o contexto, seja por gerarem blocos excessivamente extensos, com perda de precisão na busca. No presente projeto, a segmentação documental exigiu sucessivos ajustes empíricos para equilibrar granularidade semântica, preservação de contexto e restrições operacionais do banco vetorial e da camada de ingestão. Assim, a experiência prática corrobora a literatura ao indicar que o chunking não constitui etapa meramente operacional, mas um dos principais determinantes do desempenho de um sistema RAG.

Além disso, Said (2025) destaca que, em ambientes reais, a qualidade de um sistema RAG depende não apenas do pipeline básico de recuperação e geração, mas também de mecanismos de observabilidade, versionamento e avaliação. Tais mecanismos mostram-se essenciais para detectar regressões, acompanhar mudanças em modelos e parâmetros e manter consistência entre diferentes versões do sistema. Essa perspectiva foi incorporada ao projeto por meio da criação de scripts próprios de preparação, importação, avaliação e documentação dos principais parâmetros experimentais.

\subsection{Trabalhos e soluções relacionadas}

Gao et al. (2023) oferecem a base conceitual mais abrangente para compreender o RAG como resposta aos problemas de desatualização, baixa explicabilidade e geração não rastreável em LLMs. Seu trabalho orienta a decomposição do problema em etapas de recuperação, aumento de contexto e geração, fornecendo referencial teórico para o desenho da solução adotada neste estudo. No plano arquitetural, Gupta et al. (2024) consolidam o estado da arte em retrieval-augmented generation (RAG), discutindo como mecanismos de recuperação e geração se combinam para mitigar limitações de LLMs e destacando desafios de precisão, fidelidade e coordenação entre componentes.

Os trabalhos de Minaee et al. (2024) e Vaswani et al. (2017) situam a emergência dos grandes modelos de linguagem baseados em transformers e descrevem suas capacidades de compreensão e geração em ampla gama de tarefas, fornecendo o pano de fundo técnico para o uso de LLMs generativos neste projeto. Nie et al. (2024) complementam esse quadro ao tratar os text embeddings como tecnologia fundacional para tarefas de busca semântica e recuperação de informação, enfatizando a relevância da integração entre LLMs e embeddings em aplicações contemporâneas.

A literatura recente sobre alucinação em LLMs evidencia que esses modelos podem produzir saídas fluentemente redigidas, porém factualmente incorretas ou não verificáveis, o que reforça, em contextos públicos sensíveis, a necessidade de soluções explicáveis e responsáveis, com mecanismos de fundamentação e governança adequados (ANH-HOANG et al., 2025; OPENAI, 2025; HIGHBERG, 2025).

Em termos de contexto institucional, a ONU (2015), a União Interparlamentar (2024), a Westminster Foundation for Democracy (2023) e a International Idea (2023) convergem ao apontar que parlamentos e administrações públicas operam sob pressões crescentes por transparência, responsividade, responsabilização e uso qualificado da informação, o que inclui desafios específicos de organização e exploração de grandes acervos legislativos. Nesse domínio, estudos sobre discurso parlamentar e portais legislativos mostram que discursos e debates são insumos centrais para análise política e apoio técnico, mas que os sítios parlamentares frequentemente funcionam mais como repositórios informacionais do que como instrumentos efetivos de engajamento e exploração substantiva, deixando o acesso dependente de conhecimento prévio sobre estrutura e vocabulário (SKUBIC et al., 2024; BANDEIRA et al., 2016).

No campo jurídico e regulatório, Deene (2025) discute o potencial de sistemas RAG para reduzir riscos associados ao uso corporativo de inteligência artificial, sobretudo por ampliar o controle sobre as fontes, facilitar atualizações documentais e favorecer mecanismos de responsabilização. Embora o enfoque do autor esteja voltado ao contexto empresarial e à proteção de dados, seus argumentos são aplicáveis, por analogia, ao setor público, no qual a governança da informação também ocupa posição central.

No domínio político, Khaliq et al. (2024) demonstram que arquiteturas RAG podem apoiar tarefas de fact-checking político com raciocínio fundamentado em evidências externas, articulando recuperação rigorosa e geração explicável. Ainda que o objetivo deste trabalho não seja a verificação factual multimodal, a pesquisa desses autores é relevante por demonstrar a aderência do domínio político-parlamentar a abordagens baseadas em recuperação e geração com evidências.

Por fim, Martellote Viola et al. (2025) aproxima a discussão do contexto brasileiro ao tratar da inteligência artificial na organização do conhecimento legislativo. Sua contribuição reforça a pertinência de soluções explicáveis, documentalmente ancoradas e informacionalmente estruturadas em ambientes parlamentares, oferecendo base teórica específica para a aplicação proposta neste trabalho.

\subsection{Síntese da revisão e lacuna de pesquisa}

A revisão realizada permite identificar convergência entre diferentes autores quanto ao potencial do RAG para elevar a precisão factual, a auditabilidade e a governança de sistemas baseados em LLM. Verifica-se, entretanto, que parte significativa da literatura permanece concentrada em discussões conceituais amplas, cenários corporativos ou aplicações em domínios distintos do acervo legislativo brasileiro. Nesse sentido, observa-se lacuna aplicada relacionada à construção, documentação e avaliação de soluções RAG voltadas especificamente à consulta de discursos parlamentares no contexto do Senado Federal.

Assim, este trabalho busca contribuir para o preenchimento dessa lacuna ao propor, implementar e avaliar uma solução reprodutível de chatbot legislativo baseada em RAG, com foco em discursos da 56\textordfeminine{} Legislatura e atenção explícita a exigências de rastreabilidade, explicabilidade e utilidade institucional.

\section{Caracterização do problema}

\subsection{Problema da administração pública}

O problema central abordado neste trabalho consiste na dificuldade de transformar grandes acervos de dados sobre discursos legislativos em bases efetivamente consultáveis por meio de perguntas em linguagem natural, com respostas úteis, contextualizadas e verificáveis. Embora os discursos legislativos sejam públicos, seu volume, diversidade temática e dispersão temporal tornam a recuperação de argumentos, posicionamentos, temas recorrentes e referências históricas uma tarefa complexa quando realizada por mecanismos tradicionais de busca (KAVČIČ et al., 2024).

Na prática, esse quadro impacta diferentes atores institucionais. Gabinetes parlamentares podem necessitar recuperar rapidamente posicionamentos anteriores sobre determinado tema; equipes técnicas precisam localizar subsídios para notas, relatórios e pareceres; pesquisadores e jornalistas demandam acesso qualificado a discursos para fins analíticos; e cidadãos interessados em transparência e controle social podem se beneficiar de instrumentos que tornem o acervo mais acessível e inteligível (BANDEIRA et al., 2016). Na ausência de uma camada semântica de recuperação, o acesso ao conteúdo tende a ser fragmentado, lento e altamente dependente de conhecimento prévio sobre a estrutura documental e os termos empregados nos discursos (GAO et al., 2023; MARTELLOTE VIOLA et al., 2025).

Desse modo, o chatbot legislativo com RAG proposto neste trabalho procura atuar precisamente sobre essa lacuna, acrescentando uma camada de mediação informacional que combina recuperação semântica, contextualização documental e geração textual baseada em evidências.

\subsection{Partes interessadas}

As partes interessadas no projeto abrangem, em primeiro lugar, gabinetes parlamentares, consultorias legislativas e assessorias técnicas, que demandam acesso rápido e qualificado a discursos e posicionamentos anteriores. Incluem-se, também, equipes de transparência, gestão da informação e tecnologia, para as quais são particularmente relevantes aspectos de governança, manutenção, rastreabilidade e evolução da solução.

Além disso, pesquisadores das áreas de ciência política, direito, administração pública e ciência da informação constituem público potencial do sistema, assim como cidadãos e organizações da sociedade civil interessados em controle social e acompanhamento da atividade parlamentar (MARTELLOTE VIOLA et al., 2025; SAID, 2025).

\subsection{Critérios de sucesso}

Os critérios de sucesso definidos para este trabalho incluem: a capacidade de indexar uma base representativa de discursos parlamentares; a recuperação de trechos relevantes para perguntas factuais e analíticas; a geração de respostas em português fundamentadas prioritariamente no contexto recuperado; a apresentação de referências e metadados associados aos discursos utilizados; a automatização das etapas de importação e avaliação; e a reprodutibilidade da solução em ambiente local ou containerizado.

Tais critérios dialogam com recomendações da literatura para sistemas RAG em produção, que enfatizam avaliação contínua, reprodutibilidade e rastreabilidade como propriedades desejáveis em soluções voltadas a informação institucionalmente relevante (GAO et al., 2023; SAID, 2025).

\section{Proposta de solução}

\subsection{Arquitetura do chatbot legislativo}

A solução proposta materializa-se em um chatbot legislativo que utiliza a abordagem RAG para responder a consultas sobre discursos da 56\textordfeminine{} Legislatura do Senado Federal. A arquitetura compreende três serviços principais: uma instância do Open WebUI, utilizada como interface de interação e orquestração do fluxo de RAG; scripts Python empregados na extração e preparação de conteúdo documental; o ChromaDB, utilizado como banco vetorial acessado por HTTP; o Ollama para integração com grandes modelos de linguagem; e a API da OpenAI tanto para geração dos embeddings quanto para geração de respostas.

O ambiente foi definido para rodar em ambiente local usando Docker e preparado para operar tanto com embeddings locais, via Ollama, quanto com embeddings providos por serviços externos, como a API da OpenAI. Na configuração experimental validada, adotou-se esta segunda opção. Essa flexibilidade amplia a adaptabilidade da solução a diferentes restrições institucionais de infraestrutura, custo ou política tecnológica.

\subsection{Pipeline desenvolvido}

O pipeline do projeto inicia-se com a aquisição dos dados, baseada no dataset \nolinkurl{fabriciosantana/discursos-senado-legislatura-56}, que reúne discursos da 56\textordfeminine{} Legislatura do Senado Federal. Salienta-se que esse dataset foi construído em trabalho anterior a partir do portal de dados abertos do Senado. No contexto deste projeto, na etapa de pré-processamento e seleção textual, priorizou-se o campo ``TextoDiscursoIntegral'', recorrendo-se aos campos ``Resumo'' e ``Indexacao'' quando necessário, de modo a privilegiar o conteúdo mais completo e informativo na construção da base.

Em seguida, realizou-se o chunking, com geração de segmentos que preservam metadados relevantes por chunk. Esses metadados incluem data, autor, partido, unidade da federação, tipo de uso da palavra, resumo, indexação e URL do texto integral. Tal modelagem favorece não apenas a recuperação semântica, mas também a auditabilidade das respostas, uma vez que permite reconstituir a origem documental dos trechos utilizados.

A partir desses segmentos, foram gerados arquivos de ingestão para o Open WebUI em lotes estruturados em formato Markdown, além de um arquivo \texttt{jsonl} de referência. A importação dos lotes foi automatizada por script responsável por realizar o upload, monitorar o processamento e garantir a inclusão dos conteúdos na \textit{knowledge base} do Open WebUI. Concluída a ingestão, procedeu-se à indexação vetorial, com geração e persistência dos embeddings no ChromaDB.

Na etapa de consulta, o fluxo RAG integra recuperação vetorial, montagem de contexto, aplicação de prompt e geração da resposta pelo modelo. Por fim, a avaliação inicial foi realizada por meio do script \texttt{run\_rag\_eval.py}, que consulta a \textit{knowledge base} via API e produz saídas nos formatos \texttt{.jsonl}, \texttt{.md} e \texttt{.csv}, viabilizando análise estruturada, revisão qualitativa e pontuação manual para comparação entre execuções.

\subsection{Principais contribuições}

O projeto apresenta contribuições em dois níveis complementares: técnico e institucional. No plano técnico, constrói-se um ambiente reprodutível de chatbot legislativo com RAG adaptado a um acervo real do Senado Federal; demonstra-se a importância de uma modelagem de chunks enriquecida por metadados; automatizam-se etapas de importação, ingestão e avaliação; e evidencia-se empiricamente a sensibilidade do desempenho às escolhas de segmentação e conectividade com banco vetorial remoto.

No plano institucional, a solução demonstra a viabilidade de aplicar técnicas contemporâneas de recuperação e geração ao contexto do Parlamento brasileiro, aproximando recomendações gerais da literatura de RAG em produção da realidade de um acervo legislativo público. Desse modo, contribui-se para a discussão sobre governança, auditabilidade e uso responsável de inteligência artificial em contextos parlamentares.

\subsection{Riscos e limitações}

Os principais riscos e limitações identificados relacionam-se à sensibilidade do desempenho à estratégia de chunking, à possibilidade de combinação excessiva de fontes em perguntas muito abertas, à dependência da qualidade dos documentos e metadados de origem e ao risco de respostas excessivamente amplas ou inferenciais quando o prompt não é suficientemente restritivo. Em caso de limitação de acesso a recursos computacionais, há, ainda, dependência operacional de serviços externos, tanto para geração de embeddings quanto para parte da geração de respostas.

Esses riscos devem ser compreendidos em perspectiva ampliada. Como observam Deene (2025) e Said (2025), sistemas RAG podem melhorar controle e auditabilidade, mas não eliminam automaticamente problemas de qualidade documental, governança de dados, conformidade, monitoramento ou avaliação. Em outras palavras, o RAG reduz determinadas fragilidades do uso puro de LLMs, mas não substitui a necessidade de processos institucionais de revisão e controle.

\section{Experimentos e demonstração}

\subsection{Setup experimental}

O experimento utilizou uma base derivada do dataset \nolinkurl{fabriciosantana/discursos-senado-legislatura-56}, contendo 15.729 registros de entrada, dos quais 15.726 correspondiam a discursos efetivamente escritos e 3 foram descartados por ausência de texto útil. A partir desses registros, foram gerados 23.806 chunks, organizados em 120 arquivos Markdown de ingestão. Esses números foram registrados em \nolinkurl{knowledge_openwebui/build_metadata.json}, permitindo rastreabilidade quantitativa da base construída.

Na etapa de preparação da base, adotaram-se como parâmetros principais \texttt{max\_words = 850}, \texttt{overlap\_words = 150} e \texttt{chunks\_per\_file = 200}, buscando equilibrar granularidade e preservação de contexto. Já na configuração operacional validada no Open WebUI, foram definidos \textit{Chunk Size} = 6000, \textit{Chunk Overlap} = 500 e \textit{Chunk Min Size Target} = 2000, com o \textit{Markdown Header Text Splitter} habilitado, uso do modelo de embedding \texttt{text-embedding-3-small} e lote de embeddings de tamanho 32.

O prompt de RAG foi evoluído iterativamente para atender a requisitos centrais do problema institucional: privilegiar o uso prioritário do contexto recuperado, proibir explicitamente a invenção de fatos, exigir declaração de insuficiência quando a informação solicitada não estivesse presente no contexto e incentivar a distribuição de citações ao longo da resposta, complementada, quando pertinente, por uma seção final de referências. Tal configuração alinha-se à literatura que enfatiza explicabilidade, auditabilidade e governança em sistemas RAG (DEENE, 2025; SAID, 2025).

Para a avaliação, o script \texttt{run\_rag\_eval.py} realizou consultas automatizadas à \textit{knowledge base} do Open WebUI via API, gerando saídas em \texttt{.jsonl}, \texttt{.md} e \texttt{.csv}. Esse arranjo permitiu persistência estruturada, revisão qualitativa e comparação manual entre execuções, aproximando o ambiente experimental de práticas de observabilidade e versionamento recomendadas para sistemas RAG em produção.

\subsection{Resultados alcançados}

Os resultados obtidos mostraram-se satisfatórios para uma prova de conceito funcional de chatbot legislativo com RAG. A importação dos 120 lotes foi concluída com sucesso, e o sistema passou a responder perguntas factuais e temáticas sobre discursos da 56\textordfeminine{} Legislatura com qualidade inicial considerada adequada. Observou-se que perguntas específicas, delimitadas e mais atômicas apresentaram desempenho superior em comparação com perguntas muito abertas, multifocais ou excessivamente interpretativas.

A avaliação empírica também indicou que perguntas voltadas a dados objetivos ou argumentos centrais de um senador tendem a produzir respostas mais precisas, enquanto questões transversais, comparativas ou amplas elevam o risco de sínteses genéricas ou de combinação excessiva de contextos distintos. Esses achados reforçam a compreensão de que o desempenho de um sistema RAG emerge da interação entre diferentes componentes do pipeline, em especial segmentação, recuperação, seleção contextual, formulação da consulta e desenho do prompt (GAO et al., 2023; SINGH, 2025; SAID, 2025).

Verificou-se, ainda, que a inclusão de referências documentais ao final da resposta é viável por meio de instruções adequadas no prompt, o que representa ganho relevante em termos de auditabilidade. Em síntese, os resultados indicam que a solução cumpre adequadamente a função de prova de conceito, ao demonstrar viabilidade técnica, coerência arquitetural e utilidade inicial para consultas sobre o acervo discursivo analisado.

\subsection{Impacto na administração pública}

O impacto potencial do chatbot legislativo com RAG na administração pública pode ser analisado em pelo menos quatro dimensões. Em primeiro lugar, há impacto sobre a transparência, na medida em que as respostas passam a ser ancoradas em discursos públicos e citáveis, facilitando a verificação por terceiros. Em segundo lugar, há ganho de eficiência, uma vez que o sistema reduz o tempo necessário para localizar informações relevantes em grandes acervos documentais. Em terceiro lugar, a solução fortalece a memória institucional, ao facilitar a recuperação de posicionamentos anteriores, temas recorrentes e evidências documentais. Por fim, observa-se potencial de apoio à decisão, na medida em que gabinetes e equipes técnicas passam a dispor de camada adicional de mediação informacional para explorar o acervo legislativo de forma mais estruturada (MARTELLOTE VIOLA et al., 2025; DEENE, 2025).

Por operar sobre fontes controladas e permitir referência explícita aos discursos da 56\textordfeminine{} Legislatura, a solução mostra-se mais compatível com exigências de \textit{accountability} do que o uso isolado de um chatbot generalista. Essa característica assume especial relevância em contextos públicos, nos quais a confiabilidade da informação não pode ser dissociada da governança das fontes e da possibilidade de auditoria posterior.

\section{Conclusão e trabalhos futuros}

\subsection{Síntese das contribuições}

O trabalho desenvolveu uma solução funcional de chatbot legislativo com RAG voltada ao acervo de discursos da 56\textordfeminine{} Legislatura do Senado Federal. A solução compreende um pipeline de preparação documental, uma base vetorial, uma interface de consulta e scripts auxiliares de ingestão e avaliação. Demonstrou-se, assim, que é possível estruturar uma base semântica consultável e responder perguntas em linguagem natural com apoio em evidências recuperadas, aproximando o uso de LLMs das exigências de atualidade, explicabilidade e rastreabilidade apontadas pela literatura (GAO et al., 2023).

\subsection{Discussão dos resultados frente aos objetivos}

À luz dos objetivos propostos, os resultados podem ser considerados positivos. A solução atendeu aos requisitos mínimos de indexação da base de discursos da 56\textordfeminine{} Legislatura, recuperação de trechos relevantes, geração de respostas fundamentadas e automação básica da avaliação. Mais do que isso, a experiência evidenciou que a qualidade de um sistema RAG institucional depende de disciplina metodológica e de decisões arquiteturais consistentes. Estratégias de chunking, desenho do prompt, governança de metadados e avaliação iterativa influenciam diretamente o comportamento do chatbot.

Tal constatação confirma a literatura recente, segundo a qual RAG não deve ser entendido apenas como a conexão entre um LLM e um banco vetorial, mas como sistema sociotécnico que exige observabilidade, versionamento e critérios explícitos de sucesso, especialmente quando utilizado em funções públicas sensíveis (SAID, 2025). Nesse sentido, a principal contribuição do trabalho não reside apenas na implementação do protótipo, mas também na explicitação das escolhas técnicas, limitações práticas e requisitos de governança envolvidos em sua construção.

\subsection{Limitações}

Entre as limitações atuais da solução, destacam-se a avaliação ainda predominantemente qualitativa e manual; a ausência, nesta etapa, de filtros mais sofisticados por autor, período ou tipo de discurso diretamente integrados ao fluxo de geração; a possibilidade de respostas excessivamente amplas em consultas abertas; e a dependência de serviços externos para embeddings e parte da geração. Tais limitações dialogam com alertas da literatura acerca de riscos operacionais, vieses de dados e necessidade de monitoramento contínuo em sistemas RAG (GAO et al., 2023; SAID, 2025).

\subsection{Trabalhos futuros}

Como trabalhos futuros, recomenda-se ampliar a avaliação automática por meio de frameworks especializados, como RAGAS, ou de abordagens baseadas em ``juiz LLM'', de modo a complementar a avaliação humana com métricas sistemáticas. Sugere-se, também, incorporar filtros estruturados por metadado na etapa de recuperação, experimentar reranqueadores mais robustos, testar diferentes estratégias de chunking por tipo documental e implementar pós-processamento determinístico para a seção de referências.

Além disso, a solução pode ser expandida para outros tipos de documentos legislativos, como pareceres, projetos de lei, notas técnicas e documentos administrativos relacionados ao processo legislativo. Em perspectiva mais avançada, um desdobramento possível consiste na evolução do sistema para cenários de geração assistida de minutas, relatórios ou discursos, desde que preservados mecanismos de rastreabilidade, supervisão humana obrigatória e controles institucionais adequados. Dessa forma, o presente trabalho abre caminho para agenda mais ampla de pesquisa e desenvolvimento sobre inteligência artificial generativa auditável no contexto do Parlamento brasileiro.

\section*{Referências}

\small
ANH-HOANG, Dang; TRAN, Vu; NGUYEN, Le-Minh. \textit{Survey and analysis of hallucinations in large language models: attribution to prompting strategies or model behavior}. Frontiers in Artificial Intelligence, 2025. DOI: 10.3389/frai.2025.1622292. Disponível em: \url{https://www.frontiersin.org/articles/10.3389/frai.2025.1622292/full}. Acesso em: 8 abr. 2026.

BANDEIRA, Cristina Leston; BERNARDES, Cristiane Brum. \textit{Information vs engagement in parliamentary websites -- a case study of Brazil and the UK}. Revista de Sociologia e Política, Curitiba, v. 24, n. 59, p. 91--107, 2016. DOI: 10.5380/rsocp.v24i59.48709. Disponível em: \url{https://revistas.ufpr.br/rsp/article/view/48709}. Acesso em: 9 abr. 2026.

DEENE, Joris. \textit{Can a RAG system reduce the GDPR risks of your enterprise AI?}, 2025. Disponível em: \url{https://www.ictrechtswijzer.be/en/can-a-rag-system-reduce-the-gdpr-risks-of-your-corporate-ai/}. Acesso em: 9 abr. 2026.

GAO, Yunfan; XIONG, Yun; GAO, Xinyu; JIA, Kangxiang; PAN, Jinliu; BI, Yuxi; DAI, Yi; SUN, Jiawei; WANG, Meng; WANG, Haofen. \textit{Retrieval-augmented generation for large language models: a survey}. arXiv preprint, arXiv:2312.10997, 2023. Disponível em: \url{https://arxiv.org/abs/2312.10997}. Acesso em: 9 abr. 2026.

GAO, Yunfan; XIONG, Yun; WANG, Meng; WANG, Haofen. \textit{Modular RAG: Transforming RAG systems into LEGO-like reconfigurable frameworks}. arXiv preprint, arXiv:2407.21059, 2024. Disponível em: \url{https://arxiv.org/abs/2407.21059}. Acesso em: 9 abr. 2026.

GUPTA, Shailja; RANJAN, Rajesh; SINGH, Surya Narayan. \textit{A Comprehensive Survey of Retrieval-Augmented Generation (RAG): Evolution, Current Landscape and Future Directions}. arXiv preprint, arXiv:2410.12837, 2024. Disponível em: \url{https://arxiv.org/abs/2410.12837}. Acesso em: 8 abr. 2026.

HIGHBERG. \textit{Why the public sector can't do without explainable and responsible AI}. 2025. Disponível em: \url{https://highberg.com/insights/why-the-public-sector-cant-do-without-explainable-and-responsible-ai}. Acesso em: 8 abr. 2026.

INTER-PARLIAMENTARY UNION. \textit{Use cases for AI in parliaments}. Geneva, 2024. Disponível em: \url{https://www.ipu.org/file/20635/download}. Acesso em: 8 abr. 2026.

GEUNIS, Lotte. \textit{Parliamentary Monitoring Organizations: partners in parliamentary oversight}. International IDEA, 2023. Disponível em: \url{https://www.idea.int/publications/catalogue/html/parliamentary-monitoring-organizations-partners-parliamentary}. Acesso em: 8 abr. 2026.

KHALIQ, Mohammed Abdul; CHANG, Paul Yu-Chun; MA, Mingyang; PFLUGFELDER, Bernhard; MILETIĆ, Filip. \textit{RAGAR, Your Falsehood Radar: RAG-augmented reasoning for political fact-checking using multimodal large language models}. In: Proceedings of the Seventh Fact Extraction and VERification Workshop (FEVER), 2024, Miami, Florida, USA. Association for Computational Linguistics, 2024. p. 280--296. DOI: \url{https://doi.org/10.18653/v1/2024.fever-1.29}. Disponível em: \url{https://aclanthology.org/2024.fever-1.29/}. Acesso em: 9 abr. 2025.

KAVČIČ, Alenka; STOJANOSKI, Martin; MAROLT, Matija. \textit{Historical Parliamentary Corpora Viewer}. In: Proceedings of the IV Workshop on Creating, Analysing, and Increasing Accessibility of Parliamentary Corpora (ParlaCLARIN) @ LREC-COLING 2024, 2024, Torino. Disponível em: \url{https://aclanthology.org/2024.parlaclarin-1.19/}. Acesso em: 9 abr. 2026.

MAHARJAN, Anuj; YADAV, Umesh. \textit{Chunking, retrieval, and re-ranking: an empirical evaluation of RAG architectures for policy document question answering}. arXiv preprint, arXiv:2601.15457, 2026. Disponível em: \url{https://arxiv.org/abs/2601.15457}. Acesso em: 9 abr. 2026.

MARTELLOTE VIOLA, Carla Maria; MARQUES FELIPE, Carla Beatriz; FARIAS SALES MARQUES, Luana. \textit{Inteligência Artificial na Organização do Conhecimento do Parlamento brasileiro: um estudo aplicado aos dados e às informações legislativas}. ISKO Brasil, n. 8, 2025. Disponível em: \url{https://isko.org.br/ojs/index.php/iskobrasil/article/view/52}. Acesso em: 9 abr. 2026.

MINAEE, Shervin; MIKOLOV, Tomas; NIKZAD, Narjes; CHENAGHLU, Meysam; SOCHER, Richard; AMATRIAIN, Xavier; GAO, Jianfeng. \textit{Large Language Models: A Survey}. arXiv preprint, arXiv:2402.06196, 2024. Disponível em: \url{https://arxiv.org/abs/2402.06196}. Acesso em: 8 abr. 2026.

NIE, Zhijie; FENG, Zhangchi; LI, Mingxin; ZHANG, Cunwang; ZHANG, Yanzhao; LONG, Dingkun; ZHANG, Richong. \textit{When Text Embedding Meets Large Language Model: A Comprehensive Survey}. arXiv preprint, arXiv:2412.09165, 2024. Disponível em: \url{https://ar5iv.labs.arxiv.org/html/2412.09165}. Acesso em: 8 abr. 2026.

OPENAI. \textit{Why language models hallucinate}. 2025. Disponível em: \url{https://openai.com/index/why-language-models-hallucinate/}. Acesso em: 8 abr. 2026.

SAID, Adil. \textit{RAG in practice: exploring versioning, observability, and evaluation in production systems}. Towards AI, 28 ago. 2025. Disponível em: \url{https://towardsai.net/p/artificial-intelligence/rag-in-practice-exploring-versioning-observability-and-evaluation-in-production-systems}. Acesso em: 8 abr. 2026.

SINGH, Pankaj. \textit{8 types of chunking for RAG systems}. Analytics Vidhya, 4 abr. 2025. Disponível em: \url{https://www.analyticsvidhya.com/blog/2025/02/types-of-chunking-for-rag-systems/}. Acesso em: 8 abr. 2026.

SKUBIC, Jure; FIŠER, Darja. \textit{Parliamentary discourse research in political science: literature review}. In: Proceedings of the IV Workshop on Creating, Analysing, and Increasing Accessibility of Parliamentary Corpora (ParlaCLARIN) @ LREC-COLING, 2024. Torino, Italia. Disponível em: \url{https://aclanthology.org/2024.parlaclarin-1.1/}. Acesso em: 8 abr. 2026.

UNITED NATIONS. \textit{Responsive and accountable public governance}. New York: UN DESA, 2015. Disponível em: \url{https://desapublications.un.org/file/663/download}. Acesso em: 8 abr. 2026.

VASWANI, Ashish; SHAZEER, Noam; PARMAR, Niki; USZKOREIT, Jakob; JONES, Llion; GOMEZ, Aidan N.; KAISER, Lukasz; POLOSUKHIN, Illia. \textit{Attention is all you need}. arXiv preprint, arXiv:1706.03762, 2017. Disponível em: \url{https://arxiv.org/abs/1706.03762}. Acesso em: 8 abr. 2026.

WESTMINSTER FOUNDATION FOR DEMOCRACY. \textit{Guidelines for AI in Parliaments}. London, 2023. Disponível em: \url{https://www.agora-parl.org/sites/default/files/agora-documents/English%20-%20WFD%20AI%20Guidelines%20for%20Parliaments.pdf}. Acesso em: 8 abr. 2026.

\normalsize
\end{document}